\documentclass[a4paper,12pt]{article}

\usepackage[T2A]{fontenc}
\usepackage[utf8]{inputenc}
\usepackage[russian]{babel}
\usepackage{amsmath,amssymb,mathtools}
\usepackage{helvet}
\renewcommand{\familydefault}{\sfdefault}
\usepackage{icomma}
\usepackage{geometry}
\usepackage{microtype}
\usepackage{tikz}
\usepackage{tabularx,booktabs}
\usepackage{enumitem}
\usepackage{hyperref}
\usetikzlibrary{positioning}
\usepackage{parskip}
\usepackage{fancyhdr}
\usepackage{setspace}
\usepackage{xcolor}

\geometry{margin=2.1cm}
\onehalfspacing

\definecolor{accent}{HTML}{008080}
\definecolor{darkaccent}{HTML}{004D4D}
\definecolor{textgray}{HTML}{333333}
\definecolor{softgray}{HTML}{6B7474}
\definecolor{lightaccent}{HTML}{DCECEC}

\hypersetup{
    colorlinks=true,
    linkcolor=darkaccent,
    urlcolor=darkaccent
}

\setlength{\parindent}{0pt}
\setlength{\parskip}{0.6em}

\pagestyle{fancy}
\fancyhf{}
\fancyhead[L]{\small\textcolor{softgray}{Химия. Алкены}}
\fancyhead[R]{\small\textcolor{softgray}{Ксения Клыкова}}
\fancyfoot[C]{\small\textcolor{softgray}{\thepage}}
\renewcommand{\headrulewidth}{0.4pt}
\renewcommand{\headrule}{\hbox to\headwidth{\color{lightaccent}\leaders\hrule height \headrulewidth\hfill}}

\setlist[itemize]{
    leftmargin=1.6em,
    itemsep=0.3em,
    topsep=0.3em
}
\setlist[enumerate]{
    leftmargin=1.8em,
    itemsep=0.55em,
    topsep=0.4em
}

\newcommand{\topic}[1]{%
    \vspace{0.5em}
    {\Large\bfseries\textcolor{darkaccent}{#1}}\par
    \vspace{0.25em}
    {\color{accent}\rule{\linewidth}{0.7pt}}
    \vspace{0.4em}
}

\newcommand{\subtopic}[1]{%
    \vspace{0.35em}
    {\large\bfseries\textcolor{accent}{#1}}\par
}

\newcommand{\important}[1]{\textbf{\textcolor{darkaccent}{#1}}}

\begin{document}

% ---------------------------------------------------------
% ТИТУЛЬНЫЙ ЛИСТ
% ---------------------------------------------------------

\begin{titlepage}
\thispagestyle{empty}

\begin{center}

\vspace*{2.2cm}

{\Large\textcolor{softgray}{Чек-лист}}

\vspace{0.8cm}

{\fontsize{25}{30}\selectfont
\bfseries
\textcolor{darkaccent}{Алкены}}

\vspace{0.35cm}

{\Large
\textcolor{accent}{
получение, химические свойства\\
и качественные реакции
}}

\vspace{2.1cm}

{\large
\textcolor{textgray}{
Лёвин Артём Александрович\\
эксклюзивно для Ксении Клыковой
}}

\vspace{0.7cm}

{\large\textcolor{softgray}{10 класс}}

\vfill

{\large\textcolor{textgray}{\today}}

\end{center}

\vspace{1.2cm}

\begin{tikzpicture}[remember picture,overlay]
\draw[
    accent,
    line width=1.2pt
]
([xshift=1.5cm,yshift=1.4cm]current page.south west)
.. controls
([xshift=5.0cm,yshift=2.0cm]current page.south west)
and
([xshift=-5.2cm,yshift=0.75cm]current page.south east)
..
([xshift=-1.5cm,yshift=1.45cm]current page.south east);
\end{tikzpicture}

\end{titlepage}

% ---------------------------------------------------------
% ОСНОВНОЙ ЧЕК-ЛИСТ
% ---------------------------------------------------------

\topic{1. Главная структурная особенность алкенов}

\important{Алкены} --- ациклические углеводороды, содержащие одну двойную связь
\[
\mathrm{C=C}.
\]

Для алкенов с одной двойной связью общая формула:
\[
\boxed{\mathrm{C_nH_{2n}}}, \qquad n\ge2.
\]

Простейший представитель:
\[
\mathrm{CH_2=CH_2}
\]
--- этен, или этилен.

Двойная связь состоит из двух связей, одна из которых легче разрушается при химических реакциях. Поэтому для алкенов особенно характерны \important{реакции присоединения}.

Схематически:
\[
\mathrm{C=C}
\quad\longrightarrow\quad
\mathrm{C-C},
\]
а освободившиеся валентности атомов углерода используются для присоединения новых атомов или групп атомов.

\subtopic{Главная идея}

При большинстве реакций присоединения:

\begin{center}
\begin{tikzpicture}[
    node distance=1.0cm,
    every node/.style={font=\small,align=center},
    box/.style={
        draw=accent,
        rounded corners=2pt,
        inner sep=6pt,
        text width=4.3cm
    },
    arr/.style={->,>=stealth,line width=0.8pt,darkaccent}
]
\node[box] (a) {Находим двойную связь};
\node[box,right=1.3cm of a] (b) {Разрываем одну из двух связей};
\node[box,right=1.3cm of b] (c) {Присоединяем новые атомы к двум атомам углерода};

\draw[arr] (a) -- (b);
\draw[arr] (b) -- (c);
\end{tikzpicture}
\end{center}

% ---------------------------------------------------------

\topic{2. Дегидрогалогенирование: восстановление исходного вещества}

\important{Дегидрогалогенирование} --- отщепление от галогеналкана атома водорода и атома галогена с образованием двойной связи.

Типичные условия:
\[
\mathrm{KOH_{(спирт.)}},\quad t.
\]

Общая схема:
\[
\mathrm{-CH_2-CHX-}
\longrightarrow
\mathrm{-CH=CH-}
+\mathrm{HX}.
\]

В школьных уравнениях часто записывают:
\[
\mathrm{R-CH_2-CHX-R'
+KOH_{(спирт.)}
\longrightarrow
R-CH=CH-R'
+KX+H_2O}.
\]

\subtopic{Как определить возможные продукты}

Галоген отщепляется от атома углерода, при котором он расположен. Водород должен уйти от \important{соседнего} атома углерода.

Если рядом с атомом углерода, несущим галоген, имеются два разных соседних атома углерода с подходящими атомами водорода, возможно образование двух разных алкенов.

\subtopic{Пример: какой бромбутан даёт бут-1-ен и бут-2-ен?}

Рассмотрим 2-бромбутан:
\[
\mathrm{CH_3-CHBr-CH_2-CH_3}.
\]

У атома углерода с Br имеются два соседних атома углерода.

Отщепление водорода слева:
\[
\mathrm{CH_3-CHBr-CH_2-CH_3}
\longrightarrow
\mathrm{CH_2=CH-CH_2-CH_3}.
\]

Получается \important{бут-1-ен}.

Отщепление водорода справа:
\[
\mathrm{CH_3-CHBr-CH_2-CH_3}
\longrightarrow
\mathrm{CH_3-CH=CH-CH_3}.
\]

Получается \important{бут-2-ен}.

Следовательно,
\[
\boxed{\mathrm{CH_3-CHBr-CH_2-CH_3}}
\]
--- искомая структура.

Это \important{2-бромбутан}.

\subtopic{Почему 1-бромбутан не подходит?}

Для
\[
\mathrm{CH_2Br-CH_2-CH_2-CH_3}
\]
у атома углерода с Br имеется только один соседний атом углерода.

Поэтому возможное положение новой двойной связи только одно:
\[
\mathrm{CH_2=CH-CH_2-CH_3}.
\]

Получается только бут-1-ен.

% ---------------------------------------------------------

\topic{3. Правило Зайцева}

Если при дегидрогалогенировании возможно образование нескольких алкенов, преимущественно образуется более замещённый алкен.

Для 2-бромбутана:
\[
\mathrm{CH_3-CHBr-CH_2-CH_3}
\xrightarrow[\ t\ ]{\mathrm{KOH_{(спирт.)}}}
\begin{cases}
\mathrm{CH_2=CH-CH_2-CH_3},\\[2mm]
\mathrm{CH_3-CH=CH-CH_3}.
\end{cases}
\]

Основным продуктом является:
\[
\boxed{\mathrm{CH_3-CH=CH-CH_3}}
\]
--- бут-2-ен.

\important{Практический ориентир:} двойная связь преимущественно образуется там, где получается более замещённая система \(\mathrm{C=C}\).

Аналогично для 2-хлорпентана:
\[
\mathrm{CH_3-CHCl-CH_2-CH_2-CH_3}
\]
возможно образование пент-1-ена и пент-2-ена, причём преимущественно образуется \important{пент-2-ен}.

% ---------------------------------------------------------

\topic{4. Реакции присоединения к алкенам}

Для реакций присоединения полезно всегда начинать со структуры алкена и положения двойной связи.

Возьмём бут-1-ен:
\[
\mathrm{CH_2=CH-CH_2-CH_3}.
\]

\subtopic{4.1. Гидрирование}

Присоединение водорода:
\[
\mathrm{CH_2=CH-CH_2-CH_3+H_2}
\xrightarrow[\ t\ ]{\mathrm{Ni}}
\mathrm{CH_3-CH_2-CH_2-CH_3}.
\]

Продукт:
\[
\boxed{\text{бутан}}.
\]

Общая идея:
\[
\boxed{\text{алкен}+\mathrm{H_2}\longrightarrow\text{алкан}}.
\]

% ---------------------------------------------------------

\subtopic{4.2. Галогенирование}

Например, взаимодействие бут-1-ена с бромом:
\[
\mathrm{CH_2=CH-CH_2-CH_3+Br_2}
\longrightarrow
\mathrm{CH_2Br-CHBr-CH_2-CH_3}.
\]

Продукт:
\[
\boxed{\text{1,2-дибромбутан}}.
\]

Один атом Br присоединяется к одному атому углерода бывшей двойной связи, второй --- к другому.

% ---------------------------------------------------------

\subtopic{4.3. Гидратация}

\important{Гидратация} --- присоединение воды.

Для симметричного этена:
\[
\mathrm{CH_2=CH_2+H_2O}
\xrightarrow{\mathrm{H^+}}
\mathrm{CH_3-CH_2OH}.
\]

Получается этанол.

С несимметричными алкенами необходимо определить направление присоединения.

Например:
\[
\mathrm{CH_2=CH-CH_3+H_2O}
\xrightarrow{\mathrm{H^+}}
\mathrm{CH_3-CH(OH)-CH_3}.
\]

Основной продукт:
\[
\boxed{\text{пропан-2-ол}}.
\]

% ---------------------------------------------------------

\topic{5. Правило Марковникова}

При присоединении полярной молекулы типа
\[
\mathrm{H-X}
\]
к несимметричному алкену важно определить, к какому атому углерода присоединится водород.

В школьной формулировке:

\[
\boxed{
\text{H присоединяется к атому C двойной связи,}\\
\text{у которого уже больше атомов H.}
}
\]

Второй фрагмент присоединяемой молекулы направляется к другому атому углерода.

\subtopic{Пример 1. Гидратация пропена}

\[
\mathrm{CH_2=CH-CH_3+H-OH}.
\]

У первого атома углерода двойной связи имеется два атома H:
\[
\mathrm{CH_2}.
\]

У второго --- один:
\[
\mathrm{CH}.
\]

Следовательно, H преимущественно присоединяется к \(\mathrm{CH_2}\), а OH --- к \(\mathrm{CH}\):
\[
\mathrm{CH_2=CH-CH_3+H_2O}
\xrightarrow{\mathrm{H^+}}
\mathrm{CH_3-CH(OH)-CH_3}.
\]

Получается:
\[
\boxed{\text{пропан-2-ол}}.
\]

\subtopic{Пример 2. Присоединение HBr}

\[
\mathrm{CH_2=CH-CH_3+HBr}
\longrightarrow
\mathrm{CH_3-CHBr-CH_3}.
\]

Основной продукт:
\[
\boxed{\text{2-бромпропан}}.
\]

\important{Важно:} правило Марковникова используется при присоединении несимметричных реагентов к несимметричным алкенам в рассматриваемых школьных условиях.

% ---------------------------------------------------------

\topic{6. Четыре базовые реакции бут-1-ена}

Исходное вещество:
\[
\mathrm{CH_2=CH-CH_2-CH_3}.
\]

\renewcommand{\arraystretch}{1.45}
\begin{table}[ht]
\centering
\caption{Основные реакции присоединения}
\begin{tabularx}{\linewidth}{@{}lXX@{}}
\toprule
\textbf{Реагент} & \textbf{Продукт} & \textbf{Название} \\
\midrule

\(\mathrm{H_2}\) &
\(\mathrm{CH_3-CH_2-CH_2-CH_3}\) &
бутан \\

\(\mathrm{Br_2}\) &
\(\mathrm{CH_2Br-CHBr-CH_2-CH_3}\) &
1,2-дибромбутан \\

\(\mathrm{H_2O}\) &
\(\mathrm{CH_3-CH(OH)-CH_2-CH_3}\) &
бутан-2-ол \\

\(\mathrm{HBr}\) &
\(\mathrm{CH_3-CHBr-CH_2-CH_3}\) &
2-бромбутан \\

\bottomrule
\end{tabularx}
\end{table}

Для гидратации и гидробромирования направление присоединения определяется правилом Марковникова.

% ---------------------------------------------------------

\topic{7. Горение алкенов}

Алкены, как и другие углеводороды, при полном сгорании образуют:
\[
\boxed{\mathrm{CO_2+H_2O}}.
\]

Для этена:
\[
\mathrm{C_2H_4+3O_2\longrightarrow2CO_2+2H_2O}.
\]

Для пропена:
\[
\mathrm{2C_3H_6+9O_2\longrightarrow6CO_2+6H_2O}.
\]

Для бутена:
\[
\mathrm{C_4H_8+6O_2\longrightarrow4CO_2+4H_2O}.
\]

\subtopic{Как быстро уравнять горение}

Для
\[
\mathrm{C_xH_y+O_2\longrightarrow CO_2+H_2O}
\]
работайте в следующем порядке:

\begin{enumerate}
    \item Уравняйте атомы углерода.
    \item Уравняйте атомы водорода.
    \item В последнюю очередь уравняйте кислород.
    \item Если получился дробный коэффициент, умножьте все коэффициенты на подходящее целое число.
\end{enumerate}

Для алкена \(\mathrm{C_nH_{2n}}\):
\[
\mathrm{C_nH_{2n}+\frac{3n}{2}O_2
\longrightarrow nCO_2+nH_2O}.
\]

Если \(n\) нечётно, удобно умножить все коэффициенты на \(2\).

% ---------------------------------------------------------

\topic{8. Реакция Вагнера}

\important{Реакция Вагнера} --- мягкое окисление алкенов холодным разбавленным раствором перманганата калия.

Условия:
\[
\mathrm{KMnO_4},\quad \mathrm{H_2O},\quad
\text{невысокая температура}.
\]

В школьной схеме по месту двойной связи присоединяются две группы \(\mathrm{OH}\):

\[
\mathrm{CH_2=CH_2}
\xrightarrow[\mathrm{H_2O}]{\mathrm{KMnO_4}}
\mathrm{HO-CH_2-CH_2-OH}.
\]

Получается двухатомный спирт --- этан-1,2-диол.

Для бут-2-ена:
\[
\mathrm{CH_3-CH=CH-CH_3}
\xrightarrow[\mathrm{H_2O}]{\mathrm{KMnO_4}}
\mathrm{CH_3-CH(OH)-CH(OH)-CH_3}.
\]

Получается бутан-2,3-диол.

\subtopic{Наблюдение}

Раствор перманганата калия имеет интенсивную фиолетовую окраску.

При мягком окислении алкена:

\begin{itemize}
    \item фиолетовая окраска исчезает;
    \item в нейтральной или слабощелочной среде обычно образуется бурый осадок
    \(\mathrm{MnO_2}\).
\end{itemize}

Это используется как качественная проба на ненасыщенность в рамках школьного курса.

\important{Уточнение:} такая проба свидетельствует о наличии реакционноспособной ненасыщенности, однако не является абсолютно специфичной только для алкенов: с \(\mathrm{KMnO_4}\) могут реагировать и другие восстановители.

% ---------------------------------------------------------

\topic{9. Бромная вода}

\important{Бромная вода} --- водный раствор брома.

Она имеет жёлто-бурую, оранжево-бурую окраску.

Алкен присоединяет бром по двойной связи:
\[
\mathrm{CH_2=CH_2+Br_2}
\longrightarrow
\mathrm{CH_2Br-CH_2Br}.
\]

В ходе реакции окраска бромной воды исчезает.

Например:
\[
\mathrm{CH_3-CH=CH-CH_3+Br_2}
\longrightarrow
\mathrm{CH_3-CHBr-CHBr-CH_3}.
\]

Продукт:
\[
\boxed{\text{2,3-дибромбутан}}.
\]

Обесцвечивание бромной воды используют как качественную реакцию на ненасыщенность.

% ---------------------------------------------------------

\topic{10. Полимеризация}

\important{Полимеризация} --- образование крупной молекулы полимера из большого числа молекул мономера.

При полимеризации алкена двойная связь разрывается, а молекулы соединяются в длинную цепь.

\subtopic{Этен}

\[
n\,\mathrm{CH_2=CH_2}
\longrightarrow
\left(
\mathrm{-CH_2-CH_2-}
\right)_n.
\]

Мономер:
\[
\boxed{\text{этен}}.
\]

Полимер:
\[
\boxed{\text{полиэтилен}}.
\]

Структурное звено:
\[
\boxed{\mathrm{-CH_2-CH_2-}}.
\]

\subtopic{Пропен}

\[
n\,\mathrm{CH_2=CH-CH_3}
\longrightarrow
\left(
\mathrm{-CH_2-CH(CH_3)-}
\right)_n.
\]

Продукт:
\[
\boxed{\text{полипропилен}}.
\]

Структурное звено:
\[
\boxed{\mathrm{-CH_2-CH(CH_3)-}}.
\]

\subtopic{Как строить звено полимера}

\begin{enumerate}
    \item Найдите двойную связь.
    \item Замените её одинарной.
    \item Оставьте заместители при тех же атомах углерода.
    \item Покажите две связи, которыми звено соединяется с соседними звеньями.
    \item Поместите повторяющийся фрагмент в скобки и поставьте индекс \(n\).
\end{enumerate}

% ---------------------------------------------------------

\topic{11. Сводная таблица химических свойств алкенов}

\renewcommand{\arraystretch}{1.42}
\begin{table}[ht]
\centering
\caption{Что необходимо узнавать по реагенту}
\small
\begin{tabularx}{\linewidth}{@{}lXX@{}}
\toprule
\textbf{Реагент} &
\textbf{Тип реакции} &
\textbf{Главная идея} \\
\midrule

\(\mathrm{H_2}\) &
гидрирование &
образуется алкан \\

\(\mathrm{Br_2}\), \(\mathrm{Cl_2}\) &
галогенирование &
по двойной связи присоединяются два атома галогена \\

\(\mathrm{H_2O}\) &
гидратация &
образуется спирт; для несимметричных алкенов учитывают правило Марковникова \\

\(\mathrm{HCl}\), \(\mathrm{HBr}\) &
гидрогалогенирование &
образуется галогеналкан; направление определяют по правилу Марковникова \\

\(\mathrm{O_2}\) &
горение &
при полном сгорании образуются \(\mathrm{CO_2}\) и \(\mathrm{H_2O}\) \\

\(\mathrm{KMnO_4}\), \(\mathrm{H_2O}\), холод &
мягкое окисление &
образуется соединение с двумя соседними группами \(\mathrm{OH}\) \\

полимеризация &
соединение мономеров &
двойная связь раскрывается, образуется полимерная цепь \\

\bottomrule
\end{tabularx}
\end{table}

% ---------------------------------------------------------

\topic{12. Быстрый алгоритм для реакций алкенов}

\begin{center}
\begin{tikzpicture}[
    node distance=1.05cm,
    every node/.style={font=\small},
    block/.style={
        draw=darkaccent,
        rounded corners=2pt,
        inner sep=7pt,
        text width=11.7cm,
        align=left
    },
    arrow/.style={->,>=stealth,line width=0.8pt,darkaccent}
]

\node[block] (a)
{\textbf{1. Найдите двойную связь} и отметьте два атома углерода при ней.};

\node[block,below=of a] (b)
{\textbf{2. Определите тип реагента:}
\(\mathrm{H_2}\), \(\mathrm{Br_2}\), \(\mathrm{H_2O}\),
\(\mathrm{HBr}\), \(\mathrm{KMnO_4}\), \(\mathrm{O_2}\)
или условия полимеризации.};

\node[block,below=of b] (c)
{\textbf{3. Для присоединения} мысленно замените
\(\mathrm{C=C}\) на \(\mathrm{C-C}\) и распределите присоединяемые части между двумя атомами углерода.};

\node[block,below=of c] (d)
{\textbf{4. Если реагент несимметричен}, проверьте необходимость правила Марковникова.};

\node[block,below=of d] (e)
{\textbf{5. Проверьте валентность каждого атома углерода} и только затем называйте продукт.};

\draw[arrow] (a) -- (b);
\draw[arrow] (b) -- (c);
\draw[arrow] (c) -- (d);
\draw[arrow] (d) -- (e);

\end{tikzpicture}
\end{center}

% ---------------------------------------------------------

\topic{13. Типичные ошибки}

\begin{enumerate}
    \item \important{Присоединение атомов не к тем атомам углерода.}

    Новые атомы и группы при обычном присоединении появляются именно у атомов углерода, между которыми находилась двойная связь.

    \item \important{Сохранение двойной связи после присоединения.}

    В обычных реакциях присоединения рассматриваемого типа:
    \[
    \mathrm{C=C}\longrightarrow\mathrm{C-C}.
    \]

    \item \important{Игнорирование правила Марковникова.}

    Для несимметричного алкена и несимметричного реагента необходимо проверить направление присоединения.

    \item \important{Случайное ``переворачивание'' структуры в середине решения.}

    Химически зеркальная запись той же цепи допустима, однако при обучении полезнее сохранять исходное направление цепи: так легче контролировать положение групп.

    \item \important{Неправильная окраска бромной воды.}

    Бромная вода имеет жёлто-бурую или оранжево-бурую окраску. Фиолетовая окраска характерна для раствора \(\mathrm{KMnO_4}\).

    \item \important{Неправильное название спирта.}

    Используйте:
    \[
    \text{пропан-2-ол},\qquad
    \text{бутан-1-ол},\qquad
    \text{бутан-2-ол}.
    \]

    \item \important{Неверное применение правила Зайцева.}

    При отщеплении сначала определите все реально возможные положения новой двойной связи, после чего выбирайте преимущественный продукт.

    \item \important{Дробный коэффициент оставлен в окончательном уравнении.}

    Дробь допустима как промежуточный шаг. В итоговом школьном уравнении коэффициенты обычно приводят к наименьшим целым числам.
\end{enumerate}

% ---------------------------------------------------------
% ТРЕНИРОВОЧНЫЙ БЛОК
% ---------------------------------------------------------

\clearpage
\thispagestyle{fancy}

\topic{Тренировочный блок}

\textbf{Формат: Путь А.}

Выполняйте задания без подсказок. Для реакций присоединения желательно записывать развёрнутые или полусокращённые структурные формулы.

\subtopic{Средний уровень}

\begin{enumerate}

\item Запишите структурную формулу 2-бромбутана и укажите два алкена, которые могут образоваться при его дегидрогалогенировании. Какой из них является преимущественным по правилу Зайцева?

\item Завершите уравнения реакций бут-1-ена:
\[
\begin{aligned}
&\mathrm{CH_2=CH-CH_2-CH_3+H_2}
\longrightarrow ?\\
&\mathrm{CH_2=CH-CH_2-CH_3+Br_2}
\longrightarrow ?\\
&\mathrm{CH_2=CH-CH_2-CH_3+HBr}
\longrightarrow ?
\end{aligned}
\]
Назовите органические продукты.

\item Расставьте коэффициенты:
\[
\mathrm{C_5H_{10}+O_2\longrightarrow CO_2+H_2O}.
\]

\end{enumerate}

\subtopic{Выше среднего}

\begin{enumerate}[resume]

\item Для каждого вещества определите, сколько различных положений двойной связи может возникнуть при дегидрогалогенировании:

\[
\begin{aligned}
&\text{а) }\mathrm{CH_3-CH_2-CH_2Br},\\
&\text{б) }\mathrm{CH_3-CHBr-CH_2-CH_3},\\
&\text{в) }\mathrm{CH_3-CHCl-CH_2-CH_2-CH_3}.
\end{aligned}
\]

Запишите соответствующие алкены.

\item Завершите реакцию гидратации:
\[
\mathrm{CH_2=CH-CH_2-CH_3+H_2O}
\xrightarrow{\mathrm{H^+}} ?
\]
Примените правило Марковникова и назовите преимущественный продукт.

\item Завершите реакцию:
\[
\mathrm{CH_3-CH=CH_2+HCl}
\longrightarrow ?
\]
Назовите основной органический продукт и кратко укажите, как применено правило Марковникова.

\item Для бут-2-ена запишите:

\begin{enumerate}[label=\alph*)]
    \item реакцию с бромом;
    \item схему реакции Вагнера.
\end{enumerate}

Назовите оба органических продукта и укажите наблюдение в каждой качественной реакции.

\item Запишите уравнения полимеризации:

\begin{enumerate}[label=\alph*)]
    \item этена;
    \item пропена.
\end{enumerate}

Для каждого полимера выделите структурное звено.

\item Неизвестный алкен имеет формулу \(\mathrm{C_4H_8}\). При гидрировании образуется бутан, а при гидратации по правилу Марковникова преимущественно образуется бутан-2-ол. Запишите одну подходящую структурную формулу исходного алкена и уравнение его гидратации.

\end{enumerate}

\subtopic{Сложный уровень}

\begin{enumerate}[resume]

\item Галогеналкан состава \(\mathrm{C_5H_{11}Br}\) имеет неразветвлённый углеродный скелет. При дегидрогалогенировании он образует два позиционных изомера алкена, причём преимущественный продукт затем:

\begin{itemize}
    \item обесцвечивает бромную воду;
    \item при гидрировании превращается в пентан;
    \item при гидратации даёт преимущественно пентан-2-ол.
\end{itemize}

Определите структурную формулу исходного бромалкана. Запишите:

\begin{enumerate}[label=\alph*)]
    \item оба продукта дегидрогалогенирования;
    \item преимущественный продукт по правилу Зайцева;
    \item продукт его реакции с \(\mathrm{Br_2}\);
    \item продукт его гидрирования.
\end{enumerate}

\end{enumerate}

% ---------------------------------------------------------
% ОТВЕТЫ
% ---------------------------------------------------------

\clearpage
\thispagestyle{fancy}

\topic{Ответы к тренировочному блоку}

\renewcommand{\arraystretch}{1.45}

\begin{longtable}{@{}p{0.8cm}p{13.8cm}@{}}
\toprule
\textbf{№} & \textbf{Ответ} \\
\midrule
\endhead

1 &
\(\mathrm{CH_3-CHBr-CH_2-CH_3}\).

Продукты:
\[
\mathrm{CH_2=CH-CH_2-CH_3}
\quad\text{и}\quad
\mathrm{CH_3-CH=CH-CH_3}.
\]
Преимущественно образуется бут-2-ен.
\\

\midrule

2 &
\[
\mathrm{CH_2=CH-CH_2-CH_3+H_2
\rightarrow CH_3-CH_2-CH_2-CH_3}
\]
--- бутан;

\[
\mathrm{CH_2=CH-CH_2-CH_3+Br_2
\rightarrow CH_2Br-CHBr-CH_2-CH_3}
\]
--- 1,2-дибромбутан;

\[
\mathrm{CH_2=CH-CH_2-CH_3+HBr
\rightarrow CH_3-CHBr-CH_2-CH_3}
\]
--- 2-бромбутан.
\\

\midrule

3 &
\[
\boxed{
\mathrm{2C_5H_{10}+15O_2
\rightarrow10CO_2+10H_2O}
}
\]
\\

\midrule

4 &
а)
\[
\mathrm{CH_3-CH_2-CH_2Br}
\rightarrow
\mathrm{CH_3-CH=CH_2};
\]
одно положение двойной связи.

б)
\[
\mathrm{CH_3-CHBr-CH_2-CH_3}
\rightarrow
\mathrm{CH_2=CH-CH_2-CH_3}
\]
или
\[
\mathrm{CH_3-CH=CH-CH_3};
\]
два положения.

в)
\[
\mathrm{CH_3-CHCl-CH_2-CH_2-CH_3}
\rightarrow
\mathrm{CH_2=CH-CH_2-CH_2-CH_3}
\]
или
\[
\mathrm{CH_3-CH=CH-CH_2-CH_3};
\]
два положения.
\\

\midrule

5 &
\[
\mathrm{CH_2=CH-CH_2-CH_3+H_2O}
\xrightarrow{\mathrm{H^+}}
\mathrm{CH_3-CH(OH)-CH_2-CH_3}.
\]
Преимущественный продукт --- бутан-2-ол.
\\

\midrule

6 &
\[
\mathrm{CH_3-CH=CH_2+HCl
\rightarrow
CH_3-CHCl-CH_3}.
\]
Основной продукт --- 2-хлорпропан. Водород присоединяется к концевому атому углерода \(\mathrm{CH_2}\).
\\

\midrule

7 &
а)
\[
\mathrm{CH_3-CH=CH-CH_3+Br_2
\rightarrow
CH_3-CHBr-CHBr-CH_3}.
\]
Продукт --- 2,3-дибромбутан; бромная вода обесцвечивается.

б)
\[
\mathrm{CH_3-CH=CH-CH_3}
\xrightarrow[\mathrm{H_2O}]{\mathrm{KMnO_4}}
\mathrm{CH_3-CH(OH)-CH(OH)-CH_3}.
\]
Продукт --- бутан-2,3-диол; исчезает фиолетовая окраска \(\mathrm{KMnO_4}\), возможно образование бурого \(\mathrm{MnO_2}\).
\\

\midrule

8 &
а)
\[
n\,\mathrm{CH_2=CH_2}
\rightarrow
\left(\mathrm{-CH_2-CH_2-}\right)_n.
\]

б)
\[
n\,\mathrm{CH_2=CH-CH_3}
\rightarrow
\left(\mathrm{-CH_2-CH(CH_3)-}\right)_n.
\]
\\

\midrule

9 &
Один подходящий вариант --- бут-1-ен:
\[
\mathrm{CH_2=CH-CH_2-CH_3}.
\]

Гидратация:
\[
\mathrm{CH_2=CH-CH_2-CH_3+H_2O}
\xrightarrow{\mathrm{H^+}}
\mathrm{CH_3-CH(OH)-CH_2-CH_3}.
\]
Получается преимущественно бутан-2-ол.
\\

\midrule

10 &
Исходное вещество:
\[
\boxed{\mathrm{CH_3-CHBr-CH_2-CH_2-CH_3}}
\]
--- 2-бромпентан.

Дегидрогалогенирование:
\[
\mathrm{CH_3-CHBr-CH_2-CH_2-CH_3}
\rightarrow
\mathrm{CH_2=CH-CH_2-CH_2-CH_3}
\]
--- пент-1-ен,

или
\[
\mathrm{CH_3-CHBr-CH_2-CH_2-CH_3}
\rightarrow
\mathrm{CH_3-CH=CH-CH_2-CH_3}
\]
--- пент-2-ен.

По правилу Зайцева преимущественно образуется пент-2-ен.

С бромом:
\[
\mathrm{CH_3-CH=CH-CH_2-CH_3+Br_2}
\rightarrow
\mathrm{CH_3-CHBr-CHBr-CH_2-CH_3}.
\]
Продукт --- 2,3-дибромпентан.

Гидрирование:
\[
\mathrm{CH_3-CH=CH-CH_2-CH_3+H_2}
\xrightarrow{\mathrm{Ni}}
\mathrm{CH_3-CH_2-CH_2-CH_2-CH_3}.
\]
Продукт --- пентан.
\\

\bottomrule
\end{longtable}

\vfill

\begin{center}
\textcolor{darkaccent}{
\textbf{Ключевая мысль занятия:}
сначала найдите двойную связь, затем определите,
что происходит именно с двумя атомами углерода при ней.
}
\end{center}

\end{document}